\scanpage{88}
% archive foliation: 89

\uline{Proposition}\quad (i) Toute extension d'un groupe de BT par un autre est un groupe de BT.

(ii)\quad Pour tout monomorphisme de groupes de BT, le conoyau est de BT.

\deleted{(iii)\quad Pour tout épimorphisme de groupes de BT, le noyau est de BT.}

\deleted{En d'autres termes,} Si on a une suite exacte courte
\[
0\unsure{\longrightarrow|-}G'\longrightarrow G\longrightarrow G''\longrightarrow0
\]
\deleted{avec deux des termes de BT, le troisième l'est} On a une suite exacte
\[
0\longrightarrow G'(1)\unsure{-|\longrightarrow}G(1)\unsure{-|\longrightarrow}G''(1)\longrightarrow G'/pG'
\longrightarrow\bigl[G/pG\longrightarrow G''/pG''\longrightarrow0\bigr]
\]
qui dans les deux cas \unsure{enfisagés|envisagés} montre qu'on a
\[
0\longrightarrow G'(1)\unsure{-|\longrightarrow}G(1)\unsure{-|\longrightarrow}G''(1)\longrightarrow0,
\]
d'où le fait que $G(1)$ (resp. $G''(1)$) est fini loc. libre. D'autre part il est clair que $G$ (resp. $G''$) est $p$-divisible et de $p$-torsion, cqfd.

\uline{Rang ou hauteur} d'un groupe de BT. Le rang de $G(1)$ est de la forme $p^{\unsure{r|n}}$ en chaque point ($\unsure{r|n}$ une fonction localement constante sur $S$). On appelle $r$ le \uline{rang} ou la \uline{hauteur} du groupe de BT. On notera que $G(n)$ est de rang $p^{rn}$.
